\chapter{Inclusion and Exclusion}


\section{principle of Inclusion-Exclusion}

\begin{definition}
    Let $S$ be a set. Fix a collection $A_i \subset S$ indexed by $i \in \{1,\cdots,n\}$. 
    Given $I \subset \{1,\cdots,n\}$, we let $A_I= \bigcap_{i \in I}A_i$. We set $A_{\emptyset}=S$.
\end{definition}
\textbf{Remark}:
\begin{itemize}
    \item $S$ is the set of all things.
    \item $A_i$ is a set of elements in $S$ which are restricted by property indexed by $i$.
    \item $A_I$ is the set of elements which are restricted by properties $I \subset \{1,\cdots,n\}$.
    \item $A_{\emptyset}$ restricted by no property, hence is $S$.
    \item This will help us simplify our notation.
\end{itemize}

\textbf{Example}:Let $S=\{1,2,\cdots,20\}$ and let 
$$A_1=\{x \in S: x \text{ is divisible by } 2\}$$
$$A_2=\{x \in S: x \text{ is divisible by } 3\}$$
$$A_3=\{x \in S: x \text{ is divisible by } 4\}$$
\begin{itemize}
    \item $A_{\emptyset}=S=\{1,2,\cdots,20\}$
    \item $A_{\{1,2\}}=\{6,12,18\}$ these are elements in $S$ that are divisible by 6
    \item $A_{\{1,3\}}=\{4,8,12,16,20\}$ these are elements in $S$ that are divisible by 4
    \item $A_{\{1,2,3\}}=\{12\}$ these are elements in $S$ that are divisible by 12
\end{itemize}

\begin{theorem}[Basic rule]
    $$|A \cup B|=|A|+|B|-|A \cap B|$$
\end{theorem}
\begin{proof}
    We can use some basic set theory knowledge to prove this.
\end{proof}

\begin{theorem}[MIT notation]
    $$|\cup_{i=1}^nA_i|= \sum_{\emptyset \neq I \subset \{1,\cdots,n\}}(-1)^{|I|+1}|\cap_{i \in I}A_i|$$
\end{theorem}
\begin{proof}
    Induction on $n \in \N$\\
    Base case: $$|A_1|=(-1)^2|A_1|=\sum_{\emptyset \neq I \subset \{1\}}|\cap_{i \in I}A_i|$$
    Inductive step: Assume the theorem is true for $n$. Let $A_1,\cdots,A_n,A_{n+1}$ be subsets of $S$.
    \begin{equation}
        \begin{split}
            |\cup_{i=1}^{n+1}A_i| =& |\cup_{i=1}^nA_i|+|A_{n+1}|-|(\cup_{i=1}^nA_i)\cap A_{n+1}| \text{ (by the basic rule)}\\
            =& |\cup_{i=1}^nA_i|+|A_{n+1}| -|(A_1\cap A_{n+1})\cup \cdots \cup(A_n\cap A_{n+1})|\\
            =& \sum_{\emptyset \neq I \subset \{1,\cdots,n\}}(-1)^{|I|+1}|\cap_{i \in I}A_i| + |A_{n+1}|- \sum_{\emptyset \neq I \subset \{1,\cdots,n\}}(-1)^{|I|+1}|\cap_{i \in I}(A_i \cap A_{n+1})| \text{ (by I.H.)}\\
            =& \sum_{\emptyset \neq I \subset \{1,\cdots,n\}}(-1)^{|I|+1}|\cap_{i \in I}A_i| - (-1)^{|\{n+1\}|} -\sum_{J \subset \{1,\cdots,n+!\}}(-1)^{|J|}|\cap_{i \in J}A_i|\\
            & \text{ where } J \in \{I \cup \{n+1\}: \emptyset \neq I \subset \{1,\cdots,n\}\}\\
            =& \sum_{P \subset \{1,\cdots,n\}}(-1)^{|P|+1}|\cap_{i \in P}A_i| + \sum_{Q \subset \{1,\cdots,n+1\}}(-1)^{|Q|+1}|\cap_{i \in Q}A_i|\\
            & \text{ where } P \in \{I: \empty \neq I \subset \{1,\cdots,n\}\}, Q \in \{I\cup \{n+1\}: I \subset \{1,\cdots,n\}\}\\
            & \text{ Note that the first part has no $A_{n+1}$ and the second part has $A_{n+1}$}\\
            =& \sum_{\emptyset \neq T \subset \{1,\cdots,n+1\}}(-1)^{|T|+1}|\cap_{i \in T}A_i|
        \end{split}
    \end{equation}
    This is totally my proof, I think it is intuitive and direct.
\end{proof}

\begin{theorem}[Our notation]
    Let $S$ be a set. Fix a collection $A_i \subset S$ indexes by $i \in \{1,\cdots,n\}$. 
    Then we have $$|S \setminus \cup_{i=1}^nA_i|= \sum_{I \subset \{1,\cdots,n\}}(-1)^{|I|}|A_I|$$
\end{theorem}
\begin{proof}
    By induction on $n$.\\
    Base case:
    \begin{equation}
        \begin{split}
            \sum_{I \subset \{1\}}(-1)^{|I|}|A_I| &= (-1)^0|A_{\emptyset}| + (-1)^1|A_{\{1\}}|\\
            &= |S| - |A_{1}|\\
            &= |S \setminus A_1|
        \end{split}
    \end{equation}
    Inductive step: Assume the theorem is true for $n$. Let $A_1,\cdots,A_n,A_{n+1}$ be subsets of $S$. 
    Let $B_i=A_i \cap A_{n+1}$ for $i \in \{1,\cdots,n\}$ (note $B_{\emptyset}=A_{\emptyset}\cap A_{n+1}=A_{n+1}$)
    \begin{equation}
        \begin{split}
            \sum_{I \subset \{1,\cdots,n+1\}}(-1)^{|I|}|A_I| &= \sum_{I \subset \{1,\cdots,n\}}(-1)^{|I|}|A_I| + \sum_{I \subset \{1,\cdots,n\}}(-1)^{|I|+1}|A_{I \cup \{n+1\}}|\\
            &= |S \setminus \cup_{i=1}^nA_i| - \sum_{I \subset \{1,\cdots,n+1\}}(-1)^{|I|}|B_I|\\
            &= |S \setminus \cup_{i=1}^nA_i| -|B_{\emptyset} \setminus \cup_{i=1}^nB_i|\\
            &= |S \setminus \cup_{i=1}^nA_i| -|A_{n+1}\setminus \cup_{i=1}^n(A_i \cap A_{n+1})|\\
            &= |S \setminus \cup_{i=1}^nA_i| -|A_{n+1}\setminus \cup_{i=1}^nA_i|\\
            &= |(S \setminus \cup_{i=1}^nA_i) \setminus (A_{n+1} \setminus \cup_{i=1}^nA_i)|\\
            &= |(S \cup (\cup_{i=1}^nA_i)) \setminus ((\cup_{i=1}^nA_i)\cup A_{n+1})|\\
            &= |S \cup (\cup_{i=1}^{n+1}A_i)|
        \end{split}
    \end{equation}
\end{proof}
\textbf{Remark}:
\begin{itemize}
    \item We will use our notation during the course.
    \item This theorem transform counting unions to counting intersections.
    \item Intersections are easier to count since no overlap.
\end{itemize}

\noindent \textbf{Questions}:
\begin{enumerate}
    \item Compute the number of integer solutions to the equation $x_1+x_2+x_3+x_4=100$ with $x_1,x_2 \leq 10$ and $x_i \geq 0$ for all $i \in \{1,2,3,4\}$.
    \item A wizard has 5 friends. During a long wizards' conference. They met any given friend at dinner 10 times, any pair of friends 5 times, 
    any given threesome of friends 3 times, any given foursome 2 times, and all 5 friends together once. 
    The wizard ate along 6 times. Determine how many days the wizards' conference lasted.
\end{enumerate}

\noindent \textbf{Answer}:
\begin{enumerate}
    \item Let $S$ be the set of sokutions with $x_i \geq 0$ for all $i \in \{1,2,3,4\}$.\\
    Let $A_1$ be the set of solutions with $x_1 \geq 11$ and $x_i \geq 0$.\\
    Let $A_2$ be the set of solutions with $x_2 \geq 11$ and $x_i \geq 0$.\\
    We want to count $|S \setminus (A_1 \cup A_2)|$
    \begin{equation}
        \begin{split}
            |S \setminus (A_1 \cup A_2)| &= |S| -|A_1| -|A_2| + |A_1 \cap A_2|\\
            &= {100+3 \choose 3} -{89+3 \choose 3}-{89+3 \choose 3}+{78+3 \choose 3}
        \end{split}
    \end{equation}
    \item Let $S$ be the set of all days. Let $\{1,2,3,4,5\}$ denotes the set of all friends.\\
    Let $A_i=\{\text{days when the wizards met person i}\}$ for $i \in \{1,2,3,4,5\}$.\\
    From the question, we have 
    $$|A_1|=|A_2|=\cdots=|A_5|=10$$
    $$|A_{\{1,2\}}|=|A_{\{1,3\}}|=\cdots=|A_{\{4,5\}}|=5$$
    $$\forall I \subset \{1,2,3,4,5\} \text{ with } |I|=3. |A_I|=3$$
    $$\forall I \subset \{1,2,3,4,5\} \text{ with } |I|=4. |A_I|=2$$
    $$\forall I \subset \{1,2,3,4,5\} \text{ with } |I|=5. |A_I|=1$$
    $$|S \setminus \cup_{i=1}^5A_i|= |\{\text{days the wizard eats along}\}|=6$$
    By our theorem we have:
    \begin{equation}
        \begin{split}
            \sum_{I \subset \{1,\cdots,5\}}(-1)^{|I|}|A_I| &= |S| + \sum_{\emptyset \neq I \subset \{1,\cdots,n\}}(-1)^{|I|}|A_I|=6\\
            |S| &= 6 - \sum_{\emptyset \neq I \subset \{1,\cdots,n\}}(-1)^{|I|}|A_I|
        \end{split}
    \end{equation}
\end{enumerate}



\section{Counting Surjections}

\begin{definition}
    A function $f: X \longrightarrow Y$ is called a \textbf{surjection} if 
    $$\forall y \in Y. \exists x \in X. f(x)=y$$
\end{definition}

\begin{theorem}[Counting Surjections]
    Let $S(n,m)$ denotes the number of surjections $f:[n] \longrightarrow [m]$, 
    then $$S(n,m)=\sum_{k=0}^m(-1)^k{m \choose k}(m-k)^n$$
\end{theorem}
\begin{proof}
    We will use the inclusion-exclusion principle. Let $S$ be the set of all functions from $[n]$ to $[m]$. 
    For each $i \in [m]$, let $A_i=\{f: \forall x \in [n].f(x) \neq i\}$. Note, $S \setminus \cup_{i \in [n]}A_i$ is the set of surjective functions. 
    Note that for any $I \subset [m]$ of size $k$, $|A_I|=(m-k)^n$. This is because $A_I=\cap_{i \in I}A_i$ is the set of all functions that miss all of $I$.
    \begin{equation}
        \begin{split}
            S(n,m) &= |S \setminus \cup_{i \in [m]}A_i|\\
            &= \sum_{I \subset [m]}(-1)^{|I|}|A_I|\\
            &= \sum_{I \subset [m] \text{ with } |I|=0}(-1)^{|I|}|A_I| + \sum_{I \subset [m] \text{ with } |I|=1}(-1)^{|I|}|A_I| + \cdots + \sum_{I \subset [m] \text{ with } |I|=m}(-1)^{|I|}|A_I|\\
            &= \sum_{k=0}^m \sum_{I \subset [m] \text{ with } |I|=k}(-1)^{k}|A_I|\\
            &= \sum_{k=0}^m \sum_{I \subset [m] \text{ with } |I|=k}(-1)^{k}(m-k)^n\\
            &= \sum_{k=0}^m {m \choose k}(-1)^{k}|A_I|\\
        \end{split}
    \end{equation}
\end{proof}

\noindent \textbf{Question}: A teacher has 10 books(all different) that she wants to distribute to John, Paul, Ringo, George, 
ensuring that each of them gets at least one book. In how many ways can she do this? 
Compare this answer to when the books are not distinct.

\noindent \textbf{Answer}:
\begin{itemize}
    \item For different books, lets label them as $\{1,\cdots,10\}$
    $$S=\{f| f:\{1,\cdots,10\} \longrightarrow \{J,P,R,G\}\}$$
    $$|S|=S(10,4)$$
    \item For not different books, we use stars and bars:
    \begin{equation}
        \begin{cases}
            x_1+x_2+x_3+x_4 =10\\
            x_i \geq1 \text{ for all } i
        \end{cases}
    \end{equation}
    This is equivalent to:
    \begin{equation}
        \begin{cases}
            y_1+y_2+y_3+y_4=6\\
            y_i \geq 0 \text{ for all } i
        \end{cases}
    \end{equation}
\end{itemize}



\section{Counting Derangements}

\begin{definition}
    A \textbf{derangement} is a permutation $s:[N] \longrightarrow [N]$ that has no fixed points. 
    i.e. with the property that $\forall i \in[n]. s(i) \neq i$. 
    We denote the number of Derangements with $d_n$
\end{definition}

\begin{theorem}[Counting Derangements]
    $$d_n=\sum_{k=0}^n(-1)^k{n \choose k}(n-k)!=n!\sum_{k=0}^n\frac{(-1)^k}{k!}$$
\end{theorem}
\begin{proof}
    For $i \in[n]$. Let $A_i$ be the set of permutations $s$ that satisfy $s(i)=i$. 
    Let $S$ be the set of permutations of $[n]$. 
    Note that $S \setminus \cup_{i \in [n]}A_i$ is the set of all Derangements. 
    Note that for any $I \subset [n]$ of size $k$, $|A_I|=(n-k)!$. 
    Similarly to counting surjections, we have
    \begin{equation}
        \begin{split}
            d_n &= \sum_{I \subset [n]}(-1)^{|I|}|A_I|\\
        &= \sum_{k=0}^n(-1)^k{n \choose k}(n-k)!\\
        &= \sum_{k=0}^n(-1)^k \frac{n!}{k!(n-k)!}(n-k)!\\
        &= n!\sum_{k=0}^n\frac{(-1)^k}{k!}
        \end{split}
    \end{equation}
\end{proof}
\textbf{Remark}:
\begin{itemize}
    \item Consider something interesting:
    $$\frac{\text{number of derangements}}{\text{number of permutations}}=\frac{d_n}{n!}=\sum_{k=0}^n\frac{(-1)^k}{k!}$$
    We notice that this is a probability, so it is between 0 and 1. What happens when $n$ gets bigger and bigger?
    $$\sum_{k=0}^{\infty}\frac{(-1)^k}{k!}=\frac{1}{e}$$
    \item We have alternative ways to get $d_n$.
    \begin{itemize}
        \item We can use a combinatorial proof to show for $n \geq 2$, $$d_n=(n-1)(d_{n-2}+d_{n-1})$$
        \item We can do by algebra to show for $n \geq 2$, $$d_n=n \cdot d_{n-1}+(-1)^n$$
    \end{itemize}
\end{itemize}


\section{Euler's Totient Function}

\begin{definition}
    Given two integers $n,m$, we call the \textbf{greatest common divisor (GCD)} of $n$ and $m$ the largest integer which divides both $n$ and $m$. 
    We sometimes call this quantity $gcd(n,m)$.
\end{definition}

\begin{definition}
    We say two integers $n,m$ are \textbf{relatively prime} if $gcd(n,m)=1$.
\end{definition}


\begin{definition}[Euler's Totient Function]
    For $n \in \N$, define:
    $$\phi(n)=|\{m \in \{1,\cdots,n\}: gcd(m,n)=1\}|$$
\end{definition}
\textbf{Remark}:
\begin{itemize}
    \item This function is of deep interest in number theory and group theory.
\end{itemize}


\begin{lemma}
    Let $n \geq 2$ be an integer. Let $p_1,\cdots,p_k$ be a collection of distinct primes that divide $n$. 
    There are $\frac{n}{p_1\times \cdots \times p_k}$ integers in $\{1,\cdots,n\}$.
\end{lemma}
\begin{proof}
    Suppose $m \in \N$ and $m$ is divisible by all the $p_1,\cdots,p_k$. 
    It follows then that $m=p_1 \times \cdots \times p_k\times \alpha$ for some $\alpha \in \N$. So
    $$p_1 \times \cdots \times p_k \leq p_1 \times \cdots \times p_k \times \alpha \leq n$$
    This implies that $$1 \leq \alpha \leq \frac{n}{p_1 \times \cdots \times p_k}$$
\end{proof}

\begin{lemma}
    (We will use this lemma to prove the following theorem)
    $$\prod_{i=1}^k(1-x_i)=\sum_{I \subset \{1,\cdots,k\}} \prod_{i \in I}(-x_i)$$
\end{lemma}
\begin{proof}
    Induction on $k$.\\
    Base case: $k=1$ 
    \begin{equation}
        \begin{split}
            \sum_{I \subset \{1\}} \prod_{i \in I}(-x_i) &= \prod_{i \in \emptyset}(-x_i)+ \prod_{i \in \{1\}}(-x_i)\\
            &= 1 + (-x_1) \text{ (the empty product is 1)}\\
            &= \prod_{i=1}^1(1-x_i)
        \end{split}
    \end{equation}
    Inductive step:
    \begin{equation}
        \begin{split}
            \prod_{i=1}^{k+1}(1-x_i) =& (\prod_{i=1}^{k}(1-x_i))\times (1-x_{k+1})\\
            =& \sum_{I \subset \{1,\cdots,k\}} \prod_{i \in I}(-x_i) - x_{k+1} \times \sum_{I \subset \{1,\cdots,k\}} \prod_{i \in I}(-x_i)\\
            =& \sum_{P \subset \{1,\cdots,k+1\}} \prod_{i \in P}(-x_i) + \sum_{Q \subset \{1,\cdots,k+1\}} \prod_{i \in Q}(-x_i)\\
            & \text{ where } P \in \{I: I \subset \{1,\cdots,k\}\}, Q \in \{I \cup\{k+1\}: I \subset \{1,\cdots,k\}\}\\
            =& \sum_{L \subset \{1,\cdots,k+1\}} \prod_{i \in L}(-x_i)
        \end{split}
    \end{equation}
\end{proof}


\begin{theorem}[Compute $\phi(n)$]
    For all $n \geq 2$. Assume $p_1, \cdots,p_k$ are all prime factors of $n$. Then
    $$\phi(n)=n \times \prod_{i=1}^k \frac{p_i-1}{p_i}$$
\end{theorem}
\begin{proof}
    For each $i \in \{1,\cdots,k\}$. Let $A_i$ be the set of integers in $\{1,\cdots,n\}$ that are divisible by $p_i$. 
    Note, $\{1,\cdots,n\}\setminus \cup_{i=1}^kA_i$ is the set of positive integers $\leq n$ that are relatively prime with $n$. 
    By our lemma, for any $I \subset \{1,\cdots,n\}. |A_I|=|\cap_{i \in I}A_i|$. 
    Fix $i$, for $p_i$, by the lemma, there are $\frac{n}{p_i}$ numbers that are divisible by $p_i$ and $\leq n$. 
    This shows $$|A_I|=\prod_{i \in I}\frac{n}{p_i}=n \cdot \prod_{i \in I}\frac{1}{p_i}$$
    By definition of $\phi(n)$,
    \begin{equation}
        \begin{split}
            \phi(n) &= |\{1,\cdots,n\}\setminus \cup_{i=1}^kA_i|\\
            &= \sum_{I \subset \{1,\cdots,k\}}(-1)^{|I|}|A_I| \text{ (by inclusion-exclusion)}\\
            &= \sum_{I \subset \{1,\cdots,k\}}(-1)^{|I|}n \cdot \prod_{i \in I}\frac{1}{p_i}\\
            &= n \cdot \sum_{I \subset \{1,\cdots,k\}}(-1)^{|I|} \cdot \prod_{i \in I}\frac{1}{p_i}\\
            &= n \cdot \sum_{I \subset \{1,\cdots,k\}}\cdot \prod_{i \in I}\frac{(-1)}{p_i}\\
            &= n \cdot \sum_{i \in \{1,\cdots,k\}}(1-\frac{1}{p_i}) \text{ (by the lemma above)}\\
            &= n \times \prod_{i=1}^k \frac{p_i-1}{p_i}
        \end{split}
    \end{equation}
\end{proof}